\chapter{Parțiale 2018--2019}

\textbf{Numărul 1}

1. (a) Fie subspațiile vectoriale:
\begin{align*}
  V_1 &= \{ (x, y, z) \in \RR^3 \mid 2x - y + 3z = 0 \} \\
  v_2 &= \dr{Sp}\{ (1, -1, 3), (1, 2, 0) \}
\end{align*}
ale spațiului vectorial $ \RR^3 $.

Să se determine cîte o bază și dimensiunea pentru $ \RR $-spațiile vectoriale
$ V_1, V_2, V_1 + V_2, V_1 \cap V_2 $. Să se determine coordonatele vectorului
$ v = (5, 9, 13) $ după o bază a lui $ V_1 + V_2 $.

(b) Fie $ U, V \seq \RR^{20} $. Să se arate că $ V \simeq U \Leftrightarrow \dim U = \dim V $.

\vspace{0.5cm}

2. (a) Fie $ f : \RR^3 \to M_2(\RR) $ o aplicație $ \RR $-liniară, definită prin:
\[
  f(a, b, c) = %
  \begin{pmatrix}
    a + 2c & 0 \\
    a - b & b + 2c
  \end{pmatrix}.
\]

Să se determine cîte o bază și dimensiunea pentru $ \dr{Ker} f $ și $ \dr{Im}f $, matricea
asociată lui $ f $ în bazele canonice și să se determine $ V \seq \RR^3 $ astfel
încît $ \dr{Ker}f \oplus V \simeq \RR^3 $.

(b) Se consideră $ \RR^4 $ un spațiu vectorial și o aplicație liniară
$ f : \RR^4 \to \RR^4 $. Să se arate că dacă $ f $ verifică relația:
\[
  f^2 - 2f + \dr{Id}_{\RR^4} = 0_{\RR^4},
\]
atunci $ f $ este izomorfism și să se determine inversa ei. Să se determine
$ \dr{Ker}f $ și $ \dr{Im}f $ (s-a notat $ f^2 = f \circ f $).

\vspace{0.5cm}

3. (a) Să se diagonalizeze matricea
$
A = \begin{pmatrix}
  0 & 2 & 2 \\
  2 & 0 & 2 \\
  2 & 2 & 0
\end{pmatrix}
$
și apoi să se calculeze $ A^{2018} $.

(b) Fie $ A \in M_6(\RR) $ o matrice nediagonală, cu $ A^2 \neq \alpha I_6 $, pentru
orice $ \alpha \in \RR $, cu proprietatea că:
\[
  A^2 - 5 A + 6I_6 = 0_6.
\]

Să se determine valorile proprii ale matricei $ A $.

\vspace{0.5cm}
\begin{center}
  --------- $\ast\ast\ast\ast\ast$ ---------
\end{center}
\vspace{0.5cm}

\textbf{Numărul 2}

1. Fie subspațiile vectoriale:
\begin{align*}
  V_1 &= \{ (x, y, z) \in \RR^3 \mid 3x - 2y + z = 0 \} \\
  V_2 &= \dr{Sp}\{(1, 1, 2), (0, 1, 2) \}
\end{align*}
ale spațiului vectorial $ \RR^3 $.

Să se determine cîte o bază și dimensiunea pentru $ \RR $-spațiile vectoriale
$ V_1, V_2, V_1 + V_2, V_1 \cap V_2 $. Să se determine coordonatele
vectorului $ v = (15, 8, 7) $ după o bază a lui $ V_1 + V_2 $.

(b) Fie $ V $ un $ \RR $-spațiu vectorial de dimensiune 15. Să se arate
că $ V \simeq \RR^{15} $.

\vspace{0.5cm}

2. (a) Fie $ f : M_2(\RR) \to \RR^3 $ o aplicație liniară astfel încît:
\[
  f \begin{pmatrix} a & b \\ c & d \end{pmatrix} = (a - 2c + d, 0, b - c).
\]

Să se determine cîte o bază și dimensiunea pentru $ \dr{Ker} f, \dr{Im}f $,
matricea lui $ f $ în bazele canonice și să se determine $ V \seq \RR^3 $
astfel încît $ \dr{Im}f \oplus V \simeq \RR^3 $.

(b) Se consideră spațiul vectorial $ \RR^3 $ și o aplicație liniară
$ f : \RR^3 \to \RR^3 $. Să se arate că dacă $ f $ verifică relația:
\[
  f^2 - 3f - \dr{Id}_{\RR^3} = 0_{\RR^3},
\]
atunci $ f $ este izomorfism și să se determine inversa ei. Să se determine
$ \dr{Ker}f $ și $ \dr{Im}f $. Am notat cu $ f^2 = f \circ f $.

\vspace{0.5cm}

3.(a) Să se diagonalizeze matricea:
\[
  A = \begin{pmatrix}
    0 & 3 & 3 \\
    3 & 0 & 3 \\
    3 & 3 & 0
  \end{pmatrix}
\]
și apoi să se calculeze $ A^{2018} $.

(b) Fie $ f : \RR^6 \to \RR^6 $ o aplicație liniară cu proprietatea:
\[
  f^2 - 4f + 3 \dr{Id}_{\RR^6} = 0_{\RR^6}
\]
și $ f^2 \neq \alpha \dr{Id}_{\RR^6} $, pentru orice $ \alpha \in \RR $.

Să se determine valorile proprii ale matricei aplicației liniare.

%%% Local Variables:
%%% mode: latex
%%% TeX-master: "../ag-20-21"
%%% End:
